% Based on templates/lecture-example.tex; actual lecture SC3000-L10.
% Source: Week 11 Default_Logic.pptx, slides 28--41, 56--58, 64.
% Worked examples paraphrased; intermediate checks and figure reconstructed.
\begingroup
\renewcommand{\recordid}[1]{}
\renewcommand{\relatedtopic}[2]{\hyperref[topic:#1]{#2}}

\exercisegroup{SC3000-L10-EX}{第 10 讲：Default Logic}

\textbf{日期与来源：}2026-10-26；Week 11 -- Default Logic，原始 slides 28--41、56--58、64，结合视频字幕核对。以下保留讲义例题的条件与结论，补全闭包、适用性和一致性检查。

\exercisequestion{SC3000-L10-E01}{Tweety：用 Reiter 方法验证候选 extension}

\textbf{对应知识点：}\relatedtopic{SC3000-L10-T01}{默认规则}；\relatedtopic{SC3000-L10-T02}{Extension 与 Reiter 方法}。

\textbf{题目（slides 28--31）。}给定 $\Phi=\{Bird(Tweety),Cat(Sylvester)\}$ 和默认规则 $\frac{Bird(x):Fly(x)}{Fly(x)}$，检验候选
\[
 E=Cn\bigl(\Phi\cup\{Fly(Tweety)\}\bigr).
\]

\begin{workedsolution}
从 $F_0=Cn(\Phi)$ 开始。将 $x$ 实例化为 Tweety：$F_0\vdash Bird(Tweety)$，且 $E\nvdash\neg Fly(Tweety)$，因此可以加入 $Fly(Tweety)$。

对 Sylvester，$F_0$ 和后续结果均不能推出 $Bird(Sylvester)$，先决条件不成立；仅有 $Cat(Sylvester)$ 也不自动推出 $\neg Bird(Sylvester)$ 或 $\neg Fly(Sylvester)$。没有其他可增加的结论，故
\[
 F_*=Cn\bigl(\Phi\cup\{Fly(Tweety)\}\bigr)=E.
\]
候选通过验证。若后来把 $\neg Fly(Tweety)$ 加入事实，这条默认就被阻止，原有默认结论必须重新评估，不能原样保留再追加相反事实。
\end{workedsolution}

\exercisequestion{SC3000-L10-E02}{Nixon Diamond：两套扩展与两种推理标准}

\textbf{对应知识点：}\relatedtopic{SC3000-L10-T02}{多个 extension 与推理方式}。
\par\nopagebreak[4]
\textbf{题目（slides 32--41）。}在课件的玩具模型中，已知 $\Phi=\{Quaker(Nixon),Republican(Nixon)\}$，并采用
\[
 d_Q=\frac{Quaker(x):Pacifist(x)}{Pacifist(x)},\qquad
 d_R=\frac{Republican(x):\neg Pacifist(x)}{\neg Pacifist(x)}.
\]
求 extension，并判断关于 Nixon 是否为和平主义者的审慎与容许性结论。规则仅作为逻辑例题的设定，不作现实群体概括。

\begin{workedsolution}
简记 $P=Pacifist(Nixon)$。候选及其检查为：
\begin{center}
\small
\begin{tabularx}{\textwidth}{@{}lXX@{}}
\toprule
候选 & 可使用的规则 & 被阻止的规则 \\
\midrule
$E_+=Cn(\Phi\cup\{P\})$ & $d_Q$ 的前提成立，且没有 $\neg P$。 & $d_R$ 的 justification 为 $\neg P$，被候选中的 $P$ 阻止。\\
$E_-=Cn(\Phi\cup\{\neg P\})$ & $d_R$ 的前提成立，且没有 $P$。 & $d_Q$ 的 justification 为 $P$，被候选中的 $\neg P$ 阻止。\\
\bottomrule
\end{tabularx}
\end{center}
两者都能从 $\Phi$ 按 Reiter 方法构造回来，故都是 extension。若候选两种结论都不含，则两条默认都获许可，构造结果与候选不符；若候选同时含 $P,\neg P$，则两条默认均被阻止，同样无法构造回来。

$P$ 与 $\neg P$ 各自都在某个 extension 中，故分别获得 credulous 支持；没有一个出现在全部 extension 中，故两者都不是 skeptical 结论。两套 extension 必须分别考察，不能合并为 $P\land\neg P$。
\end{workedsolution}

\exercisequestion{SC3000-L10-E03}{过程顺序：成功、封闭与早先假设的冲突}

\textbf{对应知识点：}\relatedtopic{SC3000-L10-T03}{Makinson、In/Out 与过程状态}。

\textbf{题目（slides 56--58）。}给定 $\Phi=\{\alpha\}$，以及
\[
 \delta_1=\frac{\alpha:\neg\beta}{\eta},\qquad
 \delta_2=\frac{\top:\gamma}{\beta}.
\]
$\alpha,\beta,\gamma,\eta$ 为不同命题，未给出其他关系。分析 $(\delta_1)$、$(\delta_1,\delta_2)$、$(\delta_2)$，找出 extension。

\begin{workedsolution}
\textbf{先用 $\delta_1$。}当前可推出 $\alpha$，不能推出 $\beta$，因此规则可用：
\[
 In(\delta_1)=Cn(\{\alpha,\eta\}),\qquad Out(\delta_1)=\{\beta\}.
\]
二者无交集，故 successful；但 $\delta_2$ 的先决条件恒真，且没有 $\neg\gamma$，仍可应用，因此尚未 closed。

\textbf{继续用 $\delta_2$。}加入结论 $\beta$，并记录 $\neg\gamma$：
\[
 In(\delta_1,\delta_2)=Cn(\{\alpha,\eta,\beta\}),\qquad
 Out(\delta_1,\delta_2)=\{\beta,\neg\gamma\}.
\]
两条规则都已使用，过程 closed，但 $\beta\in In\cap Out$，不 successful。第二步推翻了第一步依赖的 justification $\neg\beta$。尽管 In 内没有直接出现 $\beta,\neg\beta$ 这一对矛盾，整体推理仍失败。

\textbf{改为先用 $\delta_2$。}得到
\[
 In(\delta_2)=Cn(\{\alpha,\beta\}),\qquad Out(\delta_2)=\{\neg\gamma\}.
\]
没有交集，过程 successful；已有 $\beta$ 阻止 $\delta_1$ 的 justification $\neg\beta$，因此 closed。唯一 extension 为
\[
 \boxed{E=Cn(\{\alpha,\beta\})}.
\]
结论要同时满足“没有推翻依据”和“没有漏掉仍可用的规则”。仅停在第一步的成功前缀，并不足以得到 extension。
\end{workedsolution}

\exercisequestion{SC3000-L10-E04}{过程树：局部可用的规则为何可能走向失败}

\textbf{对应知识点：}\relatedtopic{SC3000-L10-T04}{过程树}；\relatedtopic{SC3000-L10-T03}{成功与封闭}。

\textbf{题目（slide 64）。}已知事实为空，规则为
\[
 \delta_1=\frac{\top:p}{\neg q},\qquad
 \delta_2=\frac{\top:q}{r}.
\]
构造过程树，标注 In、Out 和叶节点状态，求所有 extension。

\begin{workedsolution}
根节点两条规则都可用。每个节点的第一行记录 In，第二行记录 Out；以下根据课件第 64 页重绘，并统一用 $Cn$ 表示闭包。

\begin{center}
% Original redraw of Week 11 Default_Logic.pptx, slide 64.
% Each node shows In first, Out second. Only unused applicable defaults branch.
\begin{tikzpicture}[
 >=Latex,
 state/.style={draw=noteBlue,rounded corners=2pt,align=center,
   text width=5.1cm,inner sep=6pt,font=\small},
 edge/.style={->,thick,noteBlue}
]
 \node[state,text width=3.1cm] (root) at (0,0)
   {$In=Cn(\varnothing)$\\$Out=\varnothing$};
 \node[state] (left) at (-3.25,-2.35)
   {$In=Cn(\{\neg q\})$\\$Out=\{\neg p\}$\\成功且封闭：得到 extension};
 \node[state] (right) at (3.25,-2.35)
   {$In=Cn(\{r\})$\\$Out=\{\neg q\}$\\成功，尚未封闭};
 \node[state,draw=ntuRed] (failed) at (3.25,-4.8)
   {$In=Cn(\{r,\neg q\})$\\$Out=\{\neg q,\neg p\}$\\失败：$\neg q\in In\cap Out$};
 \draw[edge] (root) -- node[midway,left,font=\small]{$\delta_1$} (left);
 \draw[edge] (root) -- node[midway,right,font=\small]{$\delta_2$} (right);
 \draw[edge] (right) -- node[right,font=\small]{$\delta_1$} (failed);
\end{tikzpicture}

\end{center}

\textbf{左支。}使用 $\delta_1$ 后，In 包含 $\neg q$，Out 为 $\{\neg p\}$，没有冲突；$\delta_2$ 的 justification $q$ 已被 $\neg q$ 阻止。该节点 successful 且 closed，给出 $Cn(\{\neg q\})$。

\textbf{右支。}使用 $\delta_2$ 后，In 为 $Cn(\{r\})$，Out 为 $\{\neg q\}$，暂时成功；但 $\delta_1$ 仍可用，尚未封闭。继续应用后推出 $\neg q$，恰好落入先前的 Out，因此失败。

两条初始分支已穷尽，唯一 extension 为 $\boxed{Cn(\{\neg q\})}$。中间节点出现过的 $r$ 只属于失败路线，不能据此声称某个 extension 支持 $r$。
\end{workedsolution}
\endgroup
